Każda z czterech analizowanych baz danych obsługuje dane półstrukturalne
w postaci dokumentów JSON, lecz robi to w sposób różniący się zarówno
implementacją, jak i konsekwencjami wydajnościowymi. W modelu danych platformy
streamingowej zostało celowo wykorzystane pole \texttt{metadata} w tabeli
\texttt{content}, przechowujące dodatkowe atrybuty marketingowe oraz
techniczne, które nie stanowią naturalnych kandydatów do reprezentacji
kolumnowej. Pole to obejmuje informacje o studiu produkcyjnym, budżecie,
otrzymanych nagrodach, znacznikach kategoryzujących, krajach koprodukcji
oraz parametrach jakościowych strumienia (maksymalna rozdzielczość, wsparcie
HDR, obsługa Dolby Atmos), które łącznie tworzą strukturę zagnieżdżoną
o zmiennej liczbie pól.

\textbf{PostgreSQL} udostępnia dwa typy danych dokumentowych: \texttt{json}
oraz \texttt{jsonb}. Pierwszy z nich zachowuje dokument w postaci tekstowej,
co skutkuje koniecznością ponownego parsowania przy każdej operacji odczytu,
natomiast \texttt{jsonb} przechowuje dokument w postaci binarnej, znormalizowanej
strukturze drzewa, eliminując narzut parsowania oraz umożliwiając efektywne
indeksowanie. W omawianym modelu pole \texttt{metadata} zostało zadeklarowane
jako \texttt{jsonb}, co pozwala na wykorzystanie pełnego zestawu operatorów
ścieżkowych: \texttt{->} (dostęp do pola jako JSON), \texttt{->>} (dostęp
do pola jako tekst), \texttt{@>} (sprawdzenie zawierania) oraz wyrażeń
\texttt{jsonpath} wprowadzonych w wersji 12. Najistotniejszą zaletą
PostgreSQL w kontekście danych półstrukturalnych pozostaje wsparcie dla
indeksów GIN, umożliwiających efektywne filtrowanie po dowolnych polach
dokumentu. Indeks GIN założony na całym dokumencie pozwala na realizację
zapytań typu „znajdź treści posiadające określony znacznik" bez konieczności
sekwencyjnego skanowania tabeli, natomiast indeks częściowy GIN wykorzystujący
wyrażenie \texttt{jsonb\_path\_ops} ogranicza rozmiar indeksu w sytuacji,
gdy wymagany jest wyłącznie operator zawierania.

\textbf{MySQL} obsługuje dane półstrukturalne poprzez natywny typ
\texttt{JSON}, dostępny od wersji 5.7. Implementacja zachowuje dokument
w wewnętrznej postaci binarnej zbliżonej do BSON, co eliminuje konieczność
parsowania tekstowego przy odczycie. Dostęp do pól realizowany jest przez
operatory \texttt{->} oraz \texttt{->>}, składniowo zgodne z PostgreSQL,
oraz przez funkcje \texttt{JSON\_EXTRACT}, \texttt{JSON\_VALUE} i
\texttt{JSON\_TABLE}. Zasadniczym ograniczeniem MySQL pozostaje brak
możliwości bezpośredniego indeksowania pól dokumentu JSON. Aby umożliwić
efektywne filtrowanie po wybranym atrybucie, należy zdefiniować kolumnę
generowaną (\texttt{GENERATED ALWAYS AS}) wyznaczającą wartość konkretnego
pola JSON, a następnie założyć na niej standardowy indeks B-tree. W modelu
omawianej platformy przykładem takiego rozwiązania jest indeks na kolumnie
generowanej wyznaczającej wartość pola \texttt{metadata.studio}, umożliwiający
szybkie filtrowanie treści według studia produkcyjnego. Rozwiązanie to,
choć funkcjonalne, jest wyraźnie mniej elastyczne niż mechanizm GIN
oferowany przez PostgreSQL, wymaga bowiem deklaratywnego określenia
indeksowanych pól na etapie projektowania schematu.

\textbf{MongoDB} traktuje dokumenty jako podstawowy model danych, co
eliminuje pojęcie dedykowanego typu JSON -- każda kolekcja przechowuje
dokumenty BSON o dowolnej strukturze. Pole \texttt{metadata} stanowi
zagnieżdżony obiekt, indeksowany w sposób analogiczny do pól zwykłych --
bezpośrednio przez ścieżkę kropkową, na przykład\\
\texttt{db.content.createIndex(\{"metadata.studio": 1\})}. Indeksy
wielokluczowe (\textit{multikey indexes}) są zakładane automatycznie dla
pól o wartości tablicowej, takich jak \texttt{metadata.tags}, co pozwala
na efektywne filtrowanie treści zawierających konkretny znacznik. Dla
zapytań po wielu polach jednocześnie dostępne są indeksy złożone, a dla
wyszukiwania pełnotekstowego w obrębie dokumentu -- indeksy tekstowe
oraz, w platformie Atlas, mechanizm Atlas Search oparty na bibliotece
Apache Lucene. W zestawieniu z silnikami relacyjnymi MongoDB oferuje
najbardziej naturalne i bezpośrednie wsparcie dla danych zagnieżdżonych,
za cenę utraty części gwarancji modelu relacyjnego, takich jak referencje
między dokumentami egzekwowane przez silnik.

\textbf{Neo4j} nie operuje typem JSON jako odrębnym typem danych --
właściwości węzłów oraz relacji są pojedynczymi wartościami skalarnymi
lub tablicowymi, lecz nie mogą posiadać dowolnej zagnieżdżonej struktury.
W konsekwencji pola, które w modelach SQL oraz dokumentowym znajdowały się
wewnątrz dokumentu \texttt{metadata}, w modelu grafowym zostały
przedstawione w postaci osobnych właściwości węzła \texttt{Content}
(\texttt{studio}, \texttt{budget}, \texttt{max\_resolution},
\texttt{hdr\_supported}, \texttt{dolby\_atmos}) lub, w przypadku list,
w postaci osobnych węzłów połączonych relacjami (na przykład węzły
\texttt{Award} mogłyby zostać połączone z \texttt{Content} relacją
\texttt{WON\_AWARD}). Biblioteka APOC udostępnia procedury
\texttt{apoc.convert.fromJsonMap} oraz \texttt{apoc.load.json}, umożliwiające
import dokumentów JSON do struktury grafowej, lecz natywny silnik Neo4j
nie przechowuje dokumentów w postaci dokumentowej.

Z porównania powyższych implementacji wynika, że PostgreSQL oraz MongoDB
oferują najbogatsze wsparcie dla danych półstrukturalnych, przy czym
PostgreSQL zachowuje dodatkowo pełen rygor relacyjny i transakcyjny,
natomiast MongoDB pozwala na operowanie dokumentami w sposób najbardziej
bezpośredni. MySQL oferuje funkcjonalność porównywalnie ograniczoną do
podstawowego dostępu i wymaga ręcznej dekompozycji indeksowanych pól,
natomiast Neo4j wymusza częściowe spłaszczenie struktury w fazie
projektowania schematu. Wybór reprezentacji pola \texttt{metadata} stanowi
zatem nie tylko decyzję modelową, lecz również istotny czynnik wpływający
na charakterystykę wydajnościową poszczególnych silników, co znalazło
odzwierciedlenie w doborze scenariuszy testowych.
